\begin{lemma}\label{lem:decic236-quotient-image}
Let $X=X_1(2,36)$ and $h(P,Q)=(P,13Q)$, and put
$Z=X/\langle h\rangle$. The curve $Z$ has genus $9$.
The image of $J_Z(\mathbb Q)$ in $J_Z(\mathbb F_5)$ is contained
in the subgroup generated by the finite cuspidal classes. This
subgroup has invariants $[3,3,3,24,24,24]$ and index $8$.
\end{lemma}
\begin{proof}
Here is an explicit model of the quotient. Over $B=X_1(2,12)$
write $(E,P,Q_{12})$ for the universal level structure, put
$T=4Q_{12}$, and let $E'=E/\langle T\rangle$.
The $3$-division polynomial of $E'$ has a linear factor for the
dual kernel. Adjoin a root of its remaining cubic factor.
This gives $Z$.
Indeed, the nine lifts $Q_{36}$ with $3Q_{36}=Q_{12}$ split into
three orbits under $h$. Their images $\phi(4Q_{36})$ give the three cyclic subgroups of
order $3$ distinct from the dual kernel.
The choice of their $x$-coordinate gives the stated cubic factor.
The model has genus $9$ and good reduction at $5$ and $7$.
Its Jacobian has rank zero, since it is a quotient of $J_X$.

Put $A_7=T_7-1-7\langle7\rangle$.
The correspondence descends from $X$, and Eichler--Shimura and
odd unramified torsion injectivity at $7$ give
$A_7J_Z(\mathbb Q)=0$.
At $5$ the proven finite Jacobian is
\[
 G=[3,3,3,24,24,192].
\]
Let $H$ be its subgroup generated by all finite cuspidal classes.
The exact computation gives $H=[3,3,3,24,24,24]$ and $G/H=C_8$.

We determine the action on these cusps by Tate labels.
For $Q_{36}=\zeta_{36}^a q^{b/36}$ and
$P=(-1)^c q^{d/2}$, impose $cb-da\equiv1\pmod2$ and quotient
by the unipotent cusp action and simultaneous negation.
The $64$ source cusps give $32$ cusps on $Z$.
Their global closed degrees are $1^4,2^{10},4^2$; at $5$ they
split as $1^8,2^{12}$.
Thirteen rational modular-unit valuation signatures distinguish
all sixteen global orbits. They identify the finite Galois
permutation $\sigma_7$ and diamond permutation $\delta_7$
unambiguously, including the four split pairs.
The boundary formula
\[
 T_7c=\sigma_7c+7\delta_7\sigma_7^{-1}c
\]
then gives $A_7H=0$ as classes.

A noncuspidal cubic place $z$ gives a class
$a=[z-3c_0]$ generating $G/H$, where $c_0$ is a global rational
cusp. Exact Riemann--Roch calculations show that $A_7a$ has order
$8$: its multiples $1,2,4$ are nonprincipal and its multiple $8$
is principal with an exported defining function.
The projected Hecke divisor has degree $24$; the kernel-orbit
algorithm and the untouched full-division-field algorithm give
identical divisors. The assignments on the twenty cusp generators
and $a$ respect every relation of $G$, independently checked by
integer lattice arithmetic. Consequently
$\ker(A_7:G\to G)=H$.
Reduction at $5$ proves the assertion.
\end{proof}

Let $D$ be an effective divisor on $X$, defined over $\mathbb Q$,
of arbitrary degree. Write its reduction at $5$ as $N+C$, with
$N$ noncuspidal and $C$ cuspidal, and let $\pi:X\to Z$ be the
quotient map. The normalized class
$[\pi_*D-(\deg D)c_0]$ lies in $J_Z(\mathbb Q)$, and hence reduces
to $H$. The normalized cuspidal class
$[\pi_*C-(\deg C)c_0]$ also lies in $H$.
It follows that
\[
 [\pi_*N-(\deg N)c_0]\in H.
\]
This necessary condition applies, in particular, when $\deg D=10$.
For a closed point of degree $d$ whose image has degree $e$, its
pushforward is $d/e$ times that image point; we include these
multiplicities in every class calculation.
